\section{Arm CCA 细粒度隔离与部署模型}

\begin{frame}[t,shrink=6]{Arm CCA 细粒度隔离与部署模型：方向开场}
\DeckKicker{方向开场}
\SlideMeta{证据等级}{方向综述}{来源状态：公开材料综合}
{\large\bfseries\color{Ink} Arm CCA 细粒度隔离与部署模型 的核心是先界定保护对象、管理边界和证据强度，再用三篇主讲材料串起技术演进。\par}\vspace{1.2mm}
\begin{columns}[T,totalwidth=\textwidth]
  \begin{column}{0.45\textwidth}
    \fcolorbox{Rule}{Paper}{%
  \begin{minipage}[t]{0.97\linewidth}
    \vspace{1.1mm}
    \hspace{1.4mm}\begin{minipage}{0.92\linewidth}
      \PanelTitle{内容说明}
      \scriptsize \begin{itemize}
  \item 从 VM 级 Realm 扩展到用户态、容器、进程内与移动/Android 部署。
  \item 固定三篇主讲材料；`liu2025lesstrust` 作为 intra-process 辅助材料，不计入三篇主讲。
  \item 先讲方向痛点和设计轴，再逐篇说明贡献、限制、证据强度和可落地条件。
  \item 对规范、Survey、draft、vendor evidence 保持显式标注，避免把背景材料写成一手系统证明。
\end{itemize}
    \end{minipage}
    \vspace{1mm}
  \end{minipage}}
  \end{column}
  \begin{column}{0.49\textwidth}
    \fcolorbox{Rule}{Panel}{%
  \begin{minipage}[t]{0.97\linewidth}
    \vspace{1.1mm}
    \hspace{1.4mm}\begin{minipage}{0.92\linewidth}
      \PanelTitle{三篇主讲材料的分工}
      \scriptsize {\tiny
\begin{tabularx}{\linewidth}{@{}YYY@{}}
\toprule
\textbf{论文} & \textbf{定位} & \textbf{证据边界} \\
\midrule
Shelter: CCA 用户态隔离 & 开创/基础入口 & E1 / 基础证据 \\
RContainer: CCA 安全容器 & 代表性改进 & E1 / 同行评审改进证据 \\
NanoZone: CCA 细粒度内存保护 & 当前 SOTA/补充边界 & E3 / 草案/未批准规范 \\
\bottomrule
\end{tabularx}
}
    \end{minipage}
    \vspace{1mm}
  \end{minipage}}
  \end{column}
\end{columns}
\SourceFootnote{来源：论文原文、官方规范或公开材料｜证据等级：见本页 badge}
\end{frame}


\begin{frame}[t,shrink=6]{Shelter: CCA 用户态隔离：内容摘要}
\DeckKicker{内容摘要}
\SlideMeta{证据等级}{E1 / 基础证据}{来源状态：本地 PDF 已核验}
{\large\bfseries\color{Ink} 用 Arm CCA 支持用户态隔离，把 Realm 能力从 VM 级机密计算推进到 user-space compartment\par}\vspace{1.2mm}
\begin{columns}[T,totalwidth=\textwidth]
  \begin{column}{0.45\textwidth}
    \fcolorbox{Rule}{Paper}{%
  \begin{minipage}[t]{0.97\linewidth}
    \vspace{1.1mm}
    \hspace{1.4mm}\begin{minipage}{0.92\linewidth}
      \PanelTitle{内容说明}
      \scriptsize \begin{itemize}
  \item 用 Arm CCA 支持用户态隔离，把 Realm 能力从 VM 级机密计算推进到 user-space compartment
  \item Shelter is the entry peer-reviewed system for moving CCA from VM Realms to user-space...
\end{itemize}
    \end{minipage}
    \vspace{1mm}
  \end{minipage}}
  \end{column}
  \begin{column}{0.49\textwidth}
    \fcolorbox{Rule}{Panel}{%
  \begin{minipage}[t]{0.97\linewidth}
    \vspace{1.1mm}
    \hspace{1.4mm}\begin{minipage}{0.92\linewidth}
      \PanelTitle{证据定位}
      \scriptsize {\tiny
\begin{tabularx}{\linewidth}{@{}YY@{}}
\toprule
\textbf{字段} & \textbf{内容} \\
\midrule
论文定位 & 开创/基础入口 \\
材料类型 & 系统论文 \\
证据等级 & E1 / 基础证据 \\
来源状态 & 本地 PDF 已核验 \\
\bottomrule
\end{tabularx}
}
    \end{minipage}
    \vspace{1mm}
  \end{minipage}}
  \end{column}
\end{columns}
\SourceFootnote{来源：Shelter: CCA 用户态隔离｜证据等级：E1 / 基础证据}
\end{frame}


\begin{frame}[t,shrink=6]{Shelter: CCA 用户态隔离：研究背景}
\DeckKicker{研究背景}
\SlideMeta{证据等级}{E1 / 基础证据}{来源状态：本地 PDF 已核验}
{\large\bfseries\color{Ink} Realm 机制不必只服务 VM\par}\vspace{1.2mm}
\begin{columns}[T,totalwidth=\textwidth]
  \begin{column}{0.45\textwidth}
    \fcolorbox{Rule}{Paper}{%
  \begin{minipage}[t]{0.97\linewidth}
    \vspace{1.1mm}
    \hspace{1.4mm}\begin{minipage}{0.92\linewidth}
      \PanelTitle{内容说明}
      \scriptsize \begin{itemize}
  \item Realm 机制不必只服务 VM
  \item 进程内敏感代码同样需要小 TCB、host-untrusted 的隔离边界
\end{itemize}
    \end{minipage}
    \vspace{1mm}
  \end{minipage}}
  \end{column}
  \begin{column}{0.49\textwidth}
    \fcolorbox{Rule}{Panel}{%
  \begin{minipage}[t]{0.97\linewidth}
    \vspace{1.1mm}
    \hspace{1.4mm}\begin{minipage}{0.92\linewidth}
      \PanelTitle{问题边界}
      \scriptsize {\tiny
\begin{tabularx}{\linewidth}{@{}YY@{}}
\toprule
\textbf{维度} & \textbf{说明} \\
\midrule
保护对象 & Arm CCA 细粒度隔离与部署模型 \\
传统瓶颈 & 从 VM 级 Realm 扩展到用户态、容器、进程内与移动/Android 部署。 \\
本文入口 & Shelter: CCA 用户态隔离 \\
证据限制 & E1 / 基础证据 \\
\bottomrule
\end{tabularx}
}
    \end{minipage}
    \vspace{1mm}
  \end{minipage}}
  \end{column}
\end{columns}
\SourceFootnote{来源：Shelter: CCA 用户态隔离｜证据等级：E1 / 基础证据}
\end{frame}


\begin{frame}[t,shrink=6]{Shelter: CCA 用户态隔离：解决方案}
\DeckKicker{解决方案}
\SlideMeta{证据等级}{E1 / 基础证据}{来源状态：本地 PDF 已核验}
{\large\bfseries\color{Ink} 把 CCA primitives 封装成用户态隔离模型，让应用 compartment 在不完全信任 OS/host 的前提下运行敏感逻辑\par}\vspace{1.2mm}
\begin{columns}[T,totalwidth=\textwidth]
  \begin{column}{0.45\textwidth}
    \fcolorbox{Rule}{Paper}{%
  \begin{minipage}[t]{0.97\linewidth}
    \vspace{1.1mm}
    \hspace{1.4mm}\begin{minipage}{0.92\linewidth}
      \PanelTitle{内容说明}
      \scriptsize \begin{itemize}
  \item 把 CCA primitives 封装成用户态隔离模型，让应用 compartment 在不完全信任 OS/host 的前提下运行敏感逻辑
\end{itemize}
    \end{minipage}
    \vspace{1mm}
  \end{minipage}}
  \end{column}
  \begin{column}{0.49\textwidth}
    \fcolorbox{Rule}{Panel}{%
  \begin{minipage}[t]{0.97\linewidth}
    \vspace{1.1mm}
    \hspace{1.4mm}\begin{minipage}{0.92\linewidth}
      \PanelTitle{核心设计拆解}
      \scriptsize \begin{itemize}
  \item 01. User-space compartment over Arm CCA
  \item 02. Realm-backed isolation boundary
  \item 03. Runtime and OS integration path for fine-grained protection
\end{itemize}
    \end{minipage}
    \vspace{1mm}
  \end{minipage}}
  \end{column}
\end{columns}
\SourceFootnote{来源：Shelter: CCA 用户态隔离｜证据等级：E1 / 基础证据}
\end{frame}


\begin{frame}[t,shrink=6]{Shelter: CCA 用户态隔离：实验结果}
\DeckKicker{实验结果}
\SlideMeta{证据等级}{E1 / 基础证据}{来源状态：本地 PDF 已核验}
{\large\bfseries\color{Ink} USENIX Security 论文提供原型与开销评估；待补证据：需要核对原文 Table/Figure 后再写入具体数值\par}\vspace{1.2mm}
\begin{columns}[T,totalwidth=\textwidth]
  \begin{column}{0.45\textwidth}
    \fcolorbox{Rule}{Paper}{%
  \begin{minipage}[t]{0.97\linewidth}
    \vspace{1.1mm}
    \hspace{1.4mm}\begin{minipage}{0.92\linewidth}
      \PanelTitle{内容说明}
      \scriptsize \begin{itemize}
  \item USENIX Security 论文提供原型与开销评估；待补证据：需要核对原文 Table/Figure 后再写入具体数值
\end{itemize}
    \end{minipage}
    \vspace{1mm}
  \end{minipage}}
  \end{column}
  \begin{column}{0.49\textwidth}
    \fcolorbox{Rule}{Panel}{%
  \begin{minipage}[t]{0.97\linewidth}
    \vspace{1.1mm}
    \hspace{1.4mm}\begin{minipage}{0.92\linewidth}
      \PanelTitle{实验/证据边界}
      \scriptsize {\tiny
\begin{tabularx}{\linewidth}{@{}YYY@{}}
\toprule
\textbf{对象} & \textbf{可支撑} & \textbf{边界} \\
\midrule
数据/来源 & Shelter: CCA 用户态隔离 & 本地 PDF 已核验 \\
Baseline/对照 & 以论文或规范原文为准 & 不跨材料外推 \\
指标/对象 & 只使用当前来源可证明的信息 & E1 / 基础证据 \\
使用边界 & 可进入本方向演进图 & 不能替代更强证据 \\
\bottomrule
\end{tabularx}
}
    \end{minipage}
    \vspace{1mm}
  \end{minipage}}
  \end{column}
\end{columns}
\SourceFootnote{来源：Shelter: CCA 用户态隔离｜证据等级：E1 / 基础证据}
\end{frame}


\begin{frame}[t,shrink=6]{Shelter: CCA 用户态隔离：文章评价}
\DeckKicker{文章评价}
\SlideMeta{证据等级}{E1 / 基础证据}{来源状态：本地 PDF 已核验}
{\large\bfseries\color{Ink} 开辟 CCA 细粒度使用方式，是从 VM Realm 走向应用隔离的方向入口\par}\vspace{1.2mm}
\begin{columns}[T,totalwidth=\textwidth]
  \begin{column}{0.45\textwidth}
    \fcolorbox{Rule}{Paper}{%
  \begin{minipage}[t]{0.97\linewidth}
    \vspace{1.1mm}
    \hspace{1.4mm}\begin{minipage}{0.92\linewidth}
      \PanelTitle{内容说明}
      \scriptsize \begin{itemize}
  \item 设计优势：开辟 CCA 细粒度使用方式，是从 VM Realm 走向应用隔离的方向入口。
  \item 局限性：系统原型不等同于通用 ABI 或量产部署，生态依赖 runtime、OS 与 CCA 硬件支持。
  \item 商业化潜力：对云端敏感函数、插件隔离和移动端安全服务有潜力；风险在开发者可用性和平台普及速度。
  \item 证据边界：E1 / Foundational；本地 PDF 已核验。
\end{itemize}
    \end{minipage}
    \vspace{1mm}
  \end{minipage}}
  \end{column}
  \begin{column}{0.49\textwidth}
    \fcolorbox{Rule}{Panel}{%
  \begin{minipage}[t]{0.97\linewidth}
    \vspace{1.1mm}
    \hspace{1.4mm}\begin{minipage}{0.92\linewidth}
      \PanelTitle{文章评价}
      \scriptsize {\tiny
\begin{tabularx}{\linewidth}{@{}YY@{}}
\toprule
\textbf{维度} & \textbf{判断} \\
\midrule
设计优势 & 开辟 CCA 细粒度使用方式，是从 VM Realm 走向应用隔离的方向入口。 \\
主要局限 & 系统原型不等同于通用 ABI 或量产部署，生态依赖 runtime、OS 与 CCA 硬件支持。 \\
商业落地 & 对云端敏感函数、插件隔离和移动端安全服务有潜力；风险在开发者可用性和平台普及速度。 \\
讲稿定位 & 开创/基础入口; 证据强度 E1 \\
\bottomrule
\end{tabularx}
}
    \end{minipage}
    \vspace{1mm}
  \end{minipage}}
  \end{column}
\end{columns}
\SourceFootnote{来源：Shelter: CCA 用户态隔离｜证据等级：E1 / 基础证据}
\end{frame}


\begin{frame}[t,shrink=6]{RContainer: CCA 安全容器：内容摘要}
\DeckKicker{内容摘要}
\SlideMeta{证据等级}{E1 / 同行评审改进证据}{来源状态：本地 PDF 已核验}
{\large\bfseries\color{Ink} 把 Arm CCA 扩展到安全容器架构，面向云原生场景重构容器隔离边界\par}\vspace{1.2mm}
\begin{columns}[T,totalwidth=\textwidth]
  \begin{column}{0.45\textwidth}
    \fcolorbox{Rule}{Paper}{%
  \begin{minipage}[t]{0.97\linewidth}
    \vspace{1.1mm}
    \hspace{1.4mm}\begin{minipage}{0.92\linewidth}
      \PanelTitle{内容说明}
      \scriptsize \begin{itemize}
  \item 把 Arm CCA 扩展到安全容器架构，面向云原生场景重构容器隔离边界
  \item RContainer is a peer-reviewed SOTA system for CCA-backed cloud-native container deployment.
\end{itemize}
    \end{minipage}
    \vspace{1mm}
  \end{minipage}}
  \end{column}
  \begin{column}{0.49\textwidth}
    \fcolorbox{Rule}{Panel}{%
  \begin{minipage}[t]{0.97\linewidth}
    \vspace{1.1mm}
    \hspace{1.4mm}\begin{minipage}{0.92\linewidth}
      \PanelTitle{证据定位}
      \scriptsize {\tiny
\begin{tabularx}{\linewidth}{@{}YY@{}}
\toprule
\textbf{字段} & \textbf{内容} \\
\midrule
论文定位 & 代表性改进 \\
材料类型 & 系统论文 \\
证据等级 & E1 / 同行评审改进证据 \\
来源状态 & 本地 PDF 已核验 \\
\bottomrule
\end{tabularx}
}
    \end{minipage}
    \vspace{1mm}
  \end{minipage}}
  \end{column}
\end{columns}
\SourceFootnote{来源：RContainer: CCA 安全容器｜证据等级：E1 / 同行评审改进证据}
\end{frame}


\begin{frame}[t,shrink=6]{RContainer: CCA 安全容器：研究背景}
\DeckKicker{研究背景}
\SlideMeta{证据等级}{E1 / 同行评审改进证据}{来源状态：本地 PDF 已核验}
{\large\bfseries\color{Ink} 容器共享内核带来部署效率，但传统容器边界难以抵御高权限 host/内核路径的泄露风险\par}\vspace{1.2mm}
\begin{columns}[T,totalwidth=\textwidth]
  \begin{column}{0.45\textwidth}
    \fcolorbox{Rule}{Paper}{%
  \begin{minipage}[t]{0.97\linewidth}
    \vspace{1.1mm}
    \hspace{1.4mm}\begin{minipage}{0.92\linewidth}
      \PanelTitle{内容说明}
      \scriptsize \begin{itemize}
  \item 容器共享内核带来部署效率，但传统容器边界难以抵御高权限 host/内核路径的泄露风险
\end{itemize}
    \end{minipage}
    \vspace{1mm}
  \end{minipage}}
  \end{column}
  \begin{column}{0.49\textwidth}
    \fcolorbox{Rule}{Panel}{%
  \begin{minipage}[t]{0.97\linewidth}
    \vspace{1.1mm}
    \hspace{1.4mm}\begin{minipage}{0.92\linewidth}
      \PanelTitle{问题边界}
      \scriptsize {\tiny
\begin{tabularx}{\linewidth}{@{}YY@{}}
\toprule
\textbf{维度} & \textbf{说明} \\
\midrule
保护对象 & Arm CCA 细粒度隔离与部署模型 \\
传统瓶颈 & 从 VM 级 Realm 扩展到用户态、容器、进程内与移动/Android 部署。 \\
本文入口 & RContainer: CCA 安全容器 \\
证据限制 & E1 / 同行评审改进证据 \\
\bottomrule
\end{tabularx}
}
    \end{minipage}
    \vspace{1mm}
  \end{minipage}}
  \end{column}
\end{columns}
\SourceFootnote{来源：RContainer: CCA 安全容器｜证据等级：E1 / 同行评审改进证据}
\end{frame}


\begin{frame}[t,shrink=6]{RContainer: CCA 安全容器：解决方案}
\DeckKicker{解决方案}
\SlideMeta{证据等级}{E1 / 同行评审改进证据}{来源状态：本地 PDF 已核验}
{\large\bfseries\color{Ink} 利用 CCA 硬件 primitives 为容器 workload 建立更强机密性边界，同时保留容器化部署模型\par}\vspace{1.2mm}
\begin{columns}[T,totalwidth=\textwidth]
  \begin{column}{0.45\textwidth}
    \fcolorbox{Rule}{Paper}{%
  \begin{minipage}[t]{0.97\linewidth}
    \vspace{1.1mm}
    \hspace{1.4mm}\begin{minipage}{0.92\linewidth}
      \PanelTitle{内容说明}
      \scriptsize \begin{itemize}
  \item 利用 CCA 硬件 primitives 为容器 workload 建立更强机密性边界，同时保留容器化部署模型
\end{itemize}
    \end{minipage}
    \vspace{1mm}
  \end{minipage}}
  \end{column}
  \begin{column}{0.49\textwidth}
    \fcolorbox{Rule}{Panel}{%
  \begin{minipage}[t]{0.97\linewidth}
    \vspace{1.1mm}
    \hspace{1.4mm}\begin{minipage}{0.92\linewidth}
      \PanelTitle{核心设计拆解}
      \scriptsize \begin{itemize}
  \item 01. CCA-backed secure container boundary
  \item 02. Container lifecycle and Realm lifecycle mapping
  \item 03. Runtime/orchestration integration considerations
\end{itemize}
    \end{minipage}
    \vspace{1mm}
  \end{minipage}}
  \end{column}
\end{columns}
\SourceFootnote{来源：RContainer: CCA 安全容器｜证据等级：E1 / 同行评审改进证据}
\end{frame}


\begin{frame}[t,shrink=6]{RContainer: CCA 安全容器：实验结果}
\DeckKicker{实验结果}
\SlideMeta{证据等级}{E1 / 同行评审改进证据}{来源状态：本地 PDF 已核验}
{\large\bfseries\color{Ink} NDSS 系统论文提供原型评估\par}\vspace{1.2mm}
\begin{columns}[T,totalwidth=\textwidth]
  \begin{column}{0.45\textwidth}
    \fcolorbox{Rule}{Paper}{%
  \begin{minipage}[t]{0.97\linewidth}
    \vspace{1.1mm}
    \hspace{1.4mm}\begin{minipage}{0.92\linewidth}
      \PanelTitle{内容说明}
      \scriptsize \begin{itemize}
  \item NDSS 系统论文提供原型评估
  \item 本报告 不复制未核对的具体性能数字
\end{itemize}
    \end{minipage}
    \vspace{1mm}
  \end{minipage}}
  \end{column}
  \begin{column}{0.49\textwidth}
    \fcolorbox{Rule}{Panel}{%
  \begin{minipage}[t]{0.97\linewidth}
    \vspace{1.1mm}
    \hspace{1.4mm}\begin{minipage}{0.92\linewidth}
      \PanelTitle{实验/证据边界}
      \scriptsize {\tiny
\begin{tabularx}{\linewidth}{@{}YYY@{}}
\toprule
\textbf{对象} & \textbf{可支撑} & \textbf{边界} \\
\midrule
数据/来源 & RContainer: CCA 安全容器 & 本地 PDF 已核验 \\
Baseline/对照 & 以论文或规范原文为准 & 不跨材料外推 \\
指标/对象 & 只使用当前来源可证明的信息 & E1 / 同行评审改进证据 \\
使用边界 & 可进入本方向演进图 & 不能替代更强证据 \\
\bottomrule
\end{tabularx}
}
    \end{minipage}
    \vspace{1mm}
  \end{minipage}}
  \end{column}
\end{columns}
\SourceFootnote{来源：RContainer: CCA 安全容器｜证据等级：E1 / 同行评审改进证据}
\end{frame}


\begin{frame}[t,shrink=6]{RContainer: CCA 安全容器：文章评价}
\DeckKicker{文章评价}
\SlideMeta{证据等级}{E1 / 同行评审改进证据}{来源状态：本地 PDF 已核验}
{\large\bfseries\color{Ink} 贴近云原生落地，把 CCA 从 confidential VM 推向容器 workload\par}\vspace{1.2mm}
\begin{columns}[T,totalwidth=\textwidth]
  \begin{column}{0.45\textwidth}
    \fcolorbox{Rule}{Paper}{%
  \begin{minipage}[t]{0.97\linewidth}
    \vspace{1.1mm}
    \hspace{1.4mm}\begin{minipage}{0.92\linewidth}
      \PanelTitle{内容说明}
      \scriptsize \begin{itemize}
  \item 设计优势：贴近云原生落地，把 CCA 从 confidential VM 推向容器 workload。
  \item 局限性：容器生态、编排系统和 CCA lifecycle 的耦合复杂，生产部署需要额外工程证据。
  \item 商业化潜力：商业潜力高，适合云平台安全容器服务；风险在 Kubernetes/runtime 改造、调试可观测性和租户迁移成本。
  \item 证据边界：E1 / Peer-reviewed SOTA；本地 PDF 已核验。
\end{itemize}
    \end{minipage}
    \vspace{1mm}
  \end{minipage}}
  \end{column}
  \begin{column}{0.49\textwidth}
    \fcolorbox{Rule}{Panel}{%
  \begin{minipage}[t]{0.97\linewidth}
    \vspace{1.1mm}
    \hspace{1.4mm}\begin{minipage}{0.92\linewidth}
      \PanelTitle{文章评价}
      \scriptsize {\tiny
\begin{tabularx}{\linewidth}{@{}YY@{}}
\toprule
\textbf{维度} & \textbf{判断} \\
\midrule
设计优势 & 贴近云原生落地，把 CCA 从 confidential VM 推向容器 workload。 \\
主要局限 & 容器生态、编排系统和 CCA lifecycle 的耦合复杂，生产部署需要额外工程证据。 \\
商业落地 & 商业潜力高，适合云平台安全容器服务；风险在 Kubernetes/runtime 改造、调试可观测性和租户迁移成本。 \\
讲稿定位 & 代表性改进; 证据强度 E1 \\
\bottomrule
\end{tabularx}
}
    \end{minipage}
    \vspace{1mm}
  \end{minipage}}
  \end{column}
\end{columns}
\SourceFootnote{来源：RContainer: CCA 安全容器｜证据等级：E1 / 同行评审改进证据}
\end{frame}


\begin{frame}[t,shrink=6]{NanoZone: CCA 细粒度内存保护：内容摘要}
\DeckKicker{内容摘要}
\SlideMeta{证据等级}{E3 / 草案/未批准规范}{来源状态：本地 PDF 已核验}
{\large\bfseries\color{Ink} 探索 Arm CCA 上更可扩展、更细粒度的内存保护域，缓解 VM/Realm 粒度过粗的问题\par}\vspace{1.2mm}
\begin{columns}[T,totalwidth=\textwidth]
  \begin{column}{0.45\textwidth}
    \fcolorbox{Rule}{Paper}{%
  \begin{minipage}[t]{0.97\linewidth}
    \vspace{1.1mm}
    \hspace{1.4mm}\begin{minipage}{0.92\linewidth}
      \PanelTitle{内容说明}
      \scriptsize \begin{itemize}
  \item 探索 Arm CCA 上更可扩展、更细粒度的内存保护域，缓解 VM/Realm 粒度过粗的问题
  \item NanoZone is retained as a draft SOTA item for fine-grained CCA memory protection, with...
\end{itemize}
    \end{minipage}
    \vspace{1mm}
  \end{minipage}}
  \end{column}
  \begin{column}{0.49\textwidth}
    \fcolorbox{Rule}{Panel}{%
  \begin{minipage}[t]{0.97\linewidth}
    \vspace{1.1mm}
    \hspace{1.4mm}\begin{minipage}{0.92\linewidth}
      \PanelTitle{证据定位}
      \scriptsize {\tiny
\begin{tabularx}{\linewidth}{@{}YY@{}}
\toprule
\textbf{字段} & \textbf{内容} \\
\midrule
论文定位 & 当前 SOTA/补充边界 \\
材料类型 & 系统论文 \\
证据等级 & E3 / 草案/未批准规范 \\
来源状态 & 本地 PDF 已核验 \\
\bottomrule
\end{tabularx}
}
    \end{minipage}
    \vspace{1mm}
  \end{minipage}}
  \end{column}
\end{columns}
\SourceFootnote{来源：NanoZone: CCA 细粒度内存保护｜证据等级：E3 / 草案/未批准规范}
\end{frame}


\begin{frame}[t,shrink=6]{NanoZone: CCA 细粒度内存保护：研究背景}
\DeckKicker{研究背景}
\SlideMeta{证据等级}{E3 / 草案/未批准规范}{来源状态：本地 PDF 已核验}
{\large\bfseries\color{Ink} 只用 Realm/VM 粒度会在管理成本、保护域数量和应用内敏感数据隔离之间形成矛盾\par}\vspace{1.2mm}
\begin{columns}[T,totalwidth=\textwidth]
  \begin{column}{0.45\textwidth}
    \fcolorbox{Rule}{Paper}{%
  \begin{minipage}[t]{0.97\linewidth}
    \vspace{1.1mm}
    \hspace{1.4mm}\begin{minipage}{0.92\linewidth}
      \PanelTitle{内容说明}
      \scriptsize \begin{itemize}
  \item 只用 Realm/VM 粒度会在管理成本、保护域数量和应用内敏感数据隔离之间形成矛盾
\end{itemize}
    \end{minipage}
    \vspace{1mm}
  \end{minipage}}
  \end{column}
  \begin{column}{0.49\textwidth}
    \fcolorbox{Rule}{Panel}{%
  \begin{minipage}[t]{0.97\linewidth}
    \vspace{1.1mm}
    \hspace{1.4mm}\begin{minipage}{0.92\linewidth}
      \PanelTitle{问题边界}
      \scriptsize {\tiny
\begin{tabularx}{\linewidth}{@{}YY@{}}
\toprule
\textbf{维度} & \textbf{说明} \\
\midrule
保护对象 & Arm CCA 细粒度隔离与部署模型 \\
传统瓶颈 & 从 VM 级 Realm 扩展到用户态、容器、进程内与移动/Android 部署。 \\
本文入口 & NanoZone: CCA 细粒度内存保护 \\
证据限制 & E3 / 草案/未批准规范 \\
\bottomrule
\end{tabularx}
}
    \end{minipage}
    \vspace{1mm}
  \end{minipage}}
  \end{column}
\end{columns}
\SourceFootnote{来源：NanoZone: CCA 细粒度内存保护｜证据等级：E3 / 草案/未批准规范}
\end{frame}


\begin{frame}[t,shrink=6]{NanoZone: CCA 细粒度内存保护：解决方案}
\DeckKicker{解决方案}
\SlideMeta{证据等级}{E3 / 草案/未批准规范}{来源状态：本地 PDF 已核验}
{\large\bfseries\color{Ink} 研究 CCA granule/lifecycle 能否支撑更小保护域，并用辅助材料 `liu2025lesstrust` 对照进程内信任收缩问题\par}\vspace{1.2mm}
\begin{columns}[T,totalwidth=\textwidth]
  \begin{column}{0.45\textwidth}
    \fcolorbox{Rule}{Paper}{%
  \begin{minipage}[t]{0.97\linewidth}
    \vspace{1.1mm}
    \hspace{1.4mm}\begin{minipage}{0.92\linewidth}
      \PanelTitle{内容说明}
      \scriptsize \begin{itemize}
  \item 研究 CCA granule/lifecycle 能否支撑更小保护域，并用辅助材料 `liu2025lesstrust` 对照进程内信任收缩问题
\end{itemize}
    \end{minipage}
    \vspace{1mm}
  \end{minipage}}
  \end{column}
  \begin{column}{0.49\textwidth}
    \fcolorbox{Rule}{Panel}{%
  \begin{minipage}[t]{0.97\linewidth}
    \vspace{1.1mm}
    \hspace{1.4mm}\begin{minipage}{0.92\linewidth}
      \PanelTitle{核心设计拆解}
      \scriptsize \begin{itemize}
  \item 01. Fine-grained memory protection over CCA
  \item 02. Granule/lifecycle scaling pressure
  \item 03. Intra-process isolation comparison with less-trust management
\end{itemize}
    \end{minipage}
    \vspace{1mm}
  \end{minipage}}
  \end{column}
\end{columns}
\SourceFootnote{来源：NanoZone: CCA 细粒度内存保护｜证据等级：E3 / 草案/未批准规范}
\end{frame}


\begin{frame}[t,shrink=6]{NanoZone: CCA 细粒度内存保护：实验结果}
\DeckKicker{实验结果}
\SlideMeta{证据等级}{E3 / 草案/未批准规范}{来源状态：本地 PDF 已核验}
{\large\bfseries\color{Ink} arXiv 预印本提供设计与评估线索\par}\vspace{1.2mm}
\begin{columns}[T,totalwidth=\textwidth]
  \begin{column}{0.45\textwidth}
    \fcolorbox{Rule}{Paper}{%
  \begin{minipage}[t]{0.97\linewidth}
    \vspace{1.1mm}
    \hspace{1.4mm}\begin{minipage}{0.92\linewidth}
      \PanelTitle{内容说明}
      \scriptsize \begin{itemize}
  \item arXiv 预印本提供设计与评估线索
  \item draft/not ratified 状态下只作为前沿方向证据，不能写成同行评议共识或成熟生产证据
\end{itemize}
    \end{minipage}
    \vspace{1mm}
  \end{minipage}}
  \end{column}
  \begin{column}{0.49\textwidth}
    \fcolorbox{Rule}{Panel}{%
  \begin{minipage}[t]{0.97\linewidth}
    \vspace{1.1mm}
    \hspace{1.4mm}\begin{minipage}{0.92\linewidth}
      \PanelTitle{实验/证据边界}
      \scriptsize {\tiny
\begin{tabularx}{\linewidth}{@{}YYY@{}}
\toprule
\textbf{对象} & \textbf{可支撑} & \textbf{边界} \\
\midrule
数据/来源 & NanoZone: CCA 细粒度内存保护 & 本地 PDF 已核验 \\
Baseline/对照 & 以论文或规范原文为准 & 不跨材料外推 \\
指标/对象 & 只使用当前来源可证明的信息 & E3 / 草案/未批准规范 \\
使用边界 & 可进入本方向演进图 & 不能替代更强证据 \\
\bottomrule
\end{tabularx}
}
    \end{minipage}
    \vspace{1mm}
  \end{minipage}}
  \end{column}
\end{columns}
\SourceFootnote{来源：NanoZone: CCA 细粒度内存保护｜证据等级：E3 / 草案/未批准规范}
\end{frame}


\begin{frame}[t,shrink=6]{NanoZone: CCA 细粒度内存保护：文章评价}
\DeckKicker{文章评价}
\SlideMeta{证据等级}{E3 / 草案/未批准规范}{来源状态：本地 PDF 已核验}
{\large\bfseries\color{Ink} 直接回应 CCA 保护粒度和可扩展性问题，适合作为细粒度隔离前沿\par}\vspace{1.2mm}
\begin{columns}[T,totalwidth=\textwidth]
  \begin{column}{0.45\textwidth}
    \fcolorbox{Rule}{Paper}{%
  \begin{minipage}[t]{0.97\linewidth}
    \vspace{1.1mm}
    \hspace{1.4mm}\begin{minipage}{0.92\linewidth}
      \PanelTitle{内容说明}
      \scriptsize \begin{itemize}
  \item 设计优势：直接回应 CCA 保护粒度和可扩展性问题，适合作为细粒度隔离前沿。
  \item 局限性：预印本尚不能代表同行评议共识或硬件/OS 生态成熟度。
  \item 商业化潜力：若硬件和 runtime 成本可控，可服务插件沙箱、密钥处理和多租户组件隔离；风险在 ABI 稳定性和开发复杂度。
  \item 证据边界：E3 / Draft-not-ratified；本地 PDF 已核验。
\end{itemize}
    \end{minipage}
    \vspace{1mm}
  \end{minipage}}
  \end{column}
  \begin{column}{0.49\textwidth}
    \fcolorbox{Rule}{Panel}{%
  \begin{minipage}[t]{0.97\linewidth}
    \vspace{1.1mm}
    \hspace{1.4mm}\begin{minipage}{0.92\linewidth}
      \PanelTitle{文章评价}
      \scriptsize {\tiny
\begin{tabularx}{\linewidth}{@{}YY@{}}
\toprule
\textbf{维度} & \textbf{判断} \\
\midrule
设计优势 & 直接回应 CCA 保护粒度和可扩展性问题，适合作为细粒度隔离前沿。 \\
主要局限 & 预印本尚不能代表同行评议共识或硬件/OS 生态成熟度。 \\
商业落地 & 若硬件和 runtime 成本可控，可服务插件沙箱、密钥处理和多租户组件隔离；风险在 ABI 稳定性和开发复杂度。 \\
讲稿定位 & 当前 SOTA/补充边界; 证据强度 E3 \\
\bottomrule
\end{tabularx}
}
    \end{minipage}
    \vspace{1mm}
  \end{minipage}}
  \end{column}
\end{columns}
\SourceFootnote{来源：NanoZone: CCA 细粒度内存保护｜证据等级：E3 / 草案/未批准规范}
\end{frame}


\begin{frame}[t,shrink=6]{Arm CCA 细粒度隔离与部署模型：方向总结}
\DeckKicker{方向总结}
\SlideMeta{证据等级}{方向综述}{来源状态：公开材料综合}
{\large\bfseries\color{Ink} Arm CCA 细粒度隔离与部署模型 的结论是：三篇主讲材料共同给出技术演进，但仍要用证据等级约束每个结论。\par}\vspace{1.2mm}
\begin{columns}[T,totalwidth=\textwidth]
  \begin{column}{0.45\textwidth}
    \fcolorbox{Rule}{Paper}{%
  \begin{minipage}[t]{0.97\linewidth}
    \vspace{1.1mm}
    \hspace{1.4mm}\begin{minipage}{0.92\linewidth}
      \PanelTitle{内容说明}
      \scriptsize \begin{itemize}
  \item Shelter: CCA 用户态隔离：开创/基础入口，E1 / 基础证据。
  \item RContainer: CCA 安全容器：代表性改进，E1 / 同行评审改进证据。
  \item NanoZone: CCA 细粒度内存保护：当前 SOTA/补充边界，E3 / 草案/未批准规范。
  \item 本方向结论：先用三篇主讲材料建立演进关系，再用证据等级控制机制、实验和商业化结论。
\end{itemize}
    \end{minipage}
    \vspace{1mm}
  \end{minipage}}
  \end{column}
  \begin{column}{0.49\textwidth}
    \fcolorbox{Rule}{Panel}{%
  \begin{minipage}[t]{0.97\linewidth}
    \vspace{1.1mm}
    \hspace{1.4mm}\begin{minipage}{0.92\linewidth}
      \PanelTitle{本方向方法对比}
      \scriptsize {\tiny
\begin{tabularx}{\linewidth}{@{}YYY@{}}
\toprule
\textbf{论文} & \textbf{定位} & \textbf{选择理由} \\
\midrule
Shelter: CCA 用户态隔离 & 开创/基础入口 & Shelter is the entry peer-reviewed system for moving CCA from VM Realms... \\
RContainer: CCA 安全容器 & 代表性改进 & RContainer is a peer-reviewed SOTA system for CCA-backed cloud-native... \\
NanoZone: CCA 细粒度内存保护 & 当前 SOTA/补充边界 & NanoZone is retained as a draft SOTA item for fine-grained CCA memory... \\
\bottomrule
\end{tabularx}
}
    \end{minipage}
    \vspace{1mm}
  \end{minipage}}
  \end{column}
\end{columns}
\SourceFootnote{来源：论文原文、官方规范或公开材料｜证据等级：见本页 badge}
\end{frame}
